% Author: Simon-Pierre Boucher — contact@spboucher.ai
%
% ═══════════════════════════════════════════════════════════════════════
% 6. ROBUSTNESS
% ═══════════════════════════════════════════════════════════════════════
\section{Robustness Checks}
\label{sec:robustness}

\subsection{Imputation Sensitivity}
\label{sec:imputation_sensitivity}

Two variables with substantial missingness---year built (19.4\%) and lot size (16.3\%)---are imputed using state-level medians. To assess the consequences, we re-estimate the baseline OLS model on subsamples that exclude imputed observations. Table~\ref{tab:imputation} reports the results.

\begin{table}[H]
\centering
\caption{Imputation Sensitivity: Unstandardized OLS on Complete-Case Subsamples}
\label{tab:imputation}
\begin{threeparttable}
\small
\begin{tabular}{lrrrr}
\toprule
Specification & $N$ & $R^2$ & $\hat{\beta}_{\ln(\text{lot})}$ & Notes \\
\midrule
Full sample (baseline) & 788,842 & 0.634 & 0.018 & All imputed values included \\
Drop missing year\_built & 640,055 & 0.638 & 0.019 & 19\% fewer obs. \\
Drop missing lot\_size & 664,900 & 0.651 & 0.059 & 16\% fewer obs. \\
Drop both missing & 534,803 & 0.654 & 0.059 & 32\% fewer obs. \\
\bottomrule
\end{tabular}
\begin{tablenotes}
\small
\item \textit{Notes:} All specifications use the same 62-regressor unstandardized OLS model with HC3 standard errors. $\hat{\beta}_{\ln(\text{lot})}$ is the coefficient on $\ln$(lot size) in natural units. Other coefficients are broadly stable across specifications (not shown for brevity).
\end{tablenotes}
\end{threeparttable}
\end{table}

The key finding is that the \textbf{lot-size coefficient triples} from 0.018 to 0.059 when imputed lot sizes are dropped. This suggests that state-median imputation attenuates the lot-size gradient substantially: imputed properties receive the state median lot size regardless of their actual lot, washing out the true lot-price relationship. The $R^2$ also increases modestly ($+2.0$ pp) when imputed lots are excluded, indicating that imputation introduces noise.

Other coefficients (living area, bathrooms, age) are broadly stable across subsamples, suggesting that the imputation concern is primarily concentrated in the lot-size variable. Dropping observations with missing year built has minimal effect on the lot-size coefficient (0.018 to 0.019) and only slightly increases $R^2$ ($+0.4$ pp).

This sensitivity finding has implications for interpretation: the baseline elasticity of listing price with respect to lot size ($0.02$) likely understates the true association, and the complete-case estimate ($0.06$) may be more informative---though it applies to a selected (non-randomly missing) subsample.


\subsection{Winsorization Sensitivity}
\label{sec:winsorization}

To assess the influence of outliers, we winsorize the two variables with implausible extreme values---lot size, capped at its 99.5th percentile (3,944 observations affected), and Bike Score, capped at its nominal maximum of 100 (38 observations)---and re-estimate the baseline OLS model.

\begin{table}[H]
\centering
\caption{Winsorization Sensitivity: Baseline vs.\ Winsorized OLS}
\label{tab:winsorization}
\begin{threeparttable}
\small
\begin{tabular}{lrr}
\toprule
Metric / Coefficient & Baseline & Winsorized \\
\midrule
$R^2$ & 0.634 & 0.640 \\
$\hat{\beta}_{\ln(\text{lot})}$ & 0.018 & 0.058 \\
$\hat{\beta}_{\ln(\text{sqft})}$ & 0.628 & 0.610 \\
$\hat{\beta}_{\text{bathrooms}}$ & 0.204 & 0.211 \\
$\hat{\beta}_{\text{garage}}$ & 0.109 & 0.111 \\
\bottomrule
\end{tabular}
\begin{tablenotes}
\small
\item \textit{Notes:} Lot size winsorized at its 99.5th percentile ($\approx 3.04$ million sqft; 3,944 observations capped); Bike Score capped at 100 (38 observations). Unstandardized OLS, $N = 788{,}842$. Other coefficients broadly stable (not shown). The large change in $\hat{\beta}_{\ln(\text{lot})}$ mirrors the imputation sensitivity finding, likely reflecting outlier lot sizes in imputed observations.
\end{tablenotes}
\end{threeparttable}
\end{table}

Winsorization increases $R^2$ modestly from 0.634 to 0.640, suggesting that extreme values in continuous variables account for a small portion of unexplained variation. The lot-size coefficient again jumps substantially (0.018 to 0.058), consistent with the imputation sensitivity analysis: extreme imputed lot sizes attenuate the gradient. Other coefficients are broadly stable, indicating that the main OLS results are not driven by outliers.


\subsection{Quantile Regression Subsample Stability}
\label{sec:qr_stability}

We repeat the median quantile regression ($\tau = 0.50$) across 10 random subsamples of 150,000 observations each. Table~\ref{tab:qr_stability} reports the coefficient of variation (CV) and sign stability for key variables.

\begin{table}[H]
\centering
\caption{Quantile Regression Stability: $\tau = 0.50$ Across 10 Subsamples}
\label{tab:qr_stability}
\begin{threeparttable}
\small
\begin{tabular}{lrrl}
\toprule
Variable & CV (\%) & Sign Stability (\%) & Assessment \\
\midrule
West & 1.0 & 100 & Highly stable \\
Northeast & 1.6 & 100 & Highly stable \\
$\ln$(Lot Size) & 2.7 & 100 & Highly stable \\
$\ln$(Living Area) & 2.9 & 100 & Highly stable \\
Bathrooms & 2.9 & 100 & Highly stable \\
Garage & 4.1 & 100 & Highly stable \\
Luxury Score & 4.8 & 100 & Highly stable \\
Age$^2$ & 5.3 & 100 & Highly stable \\
Foreclosure & 7.1 & 100 & Stable \\
Bedrooms & 7.3 & 100 & Stable \\
\midrule
\multicolumn{4}{l}{\textit{Less stable variables}} \\
Pool & 33.4 & 100 & Moderately unstable \\
Waterfront & 42.2 & 100 & Moderately unstable \\
Age & 78.6 & 80 & Unstable \\
\bottomrule
\end{tabular}
\begin{tablenotes}
\small
\item \textit{Notes:} CV = coefficient of variation across 10 random subsamples of 150,000 observations. Sign stability = percentage of subsamples in which the coefficient has the same sign as the pooled estimate. Variables sorted by CV. The instability of age (CV = 78.6\%) likely reflects heterogeneity in the age-price relationship across geographic subsamples. The waterfront and pool coefficients vary in magnitude but never change sign; their instability reflects low variation in these binary variables in some subsamples.
\end{tablenotes}
\end{threeparttable}
\end{table}

Most variables exhibit high stability (CV $< 8\%$, 100\% sign stability), confirming that the quantile regression results are robust to subsample variation. Three variables are exceptions: \textbf{age} (CV = 78.6\%, sign stability = 80\%), \textbf{waterfront} (CV = 42.2\%), and \textbf{pool} (CV = 33.4\%). The instability of age is consistent with substantial geographic heterogeneity in the age-price relationship---old properties may command a premium in some markets (historic charm) and a discount in others (obsolescence). Waterfront and pool instability reflects the relatively low prevalence of these features in some subsamples.


\subsection{Robustness Summary Matrix}
\label{sec:robustness_matrix}

Table~\ref{tab:robustness_matrix} provides a comprehensive summary of all robustness findings.

\begin{table}[H]
\centering
\caption{Robustness Summary Matrix}
\label{tab:robustness_matrix}
\begin{threeparttable}
\footnotesize
\begin{tabular}{p{3.5cm}p{2.5cm}p{2.5cm}p{4.5cm}}
\toprule
Test & Baseline & Robustness & Key Finding \\
\midrule
\multicolumn{4}{l}{\textit{OLS Specification}} \\
Region FE $\rightarrow$ State FE & $R^2 = 0.630$ & $R^2 = 0.678$ & $+$4.8 pp from state-level controls \\
Region FE $\rightarrow$ ZIP3 FE & $R^2 = 0.630$ & $R^2 = 0.725$ & $+$9.5 pp from ZIP3 controls \\
Drop imputed lot sizes & $\hat{\beta}_{\ln\text{lot}} = 0.018$ & $\hat{\beta}_{\ln\text{lot}} = 0.059$ & Coefficient triples; imputation attenuates \\
Winsorization (1\%/99\%) & $R^2 = 0.634$ & $R^2 = 0.640$ & Modest improvement; lot coeff.\ sensitive \\
\midrule
\multicolumn{4}{l}{\textit{Spatial Diagnostics}} \\
Moran's $I$ (3 subsamples) & $I = 0.2745$ & std $= 0.0076$ & Highly stable; strong spatial autocorrelation \\
\midrule
\multicolumn{4}{l}{\textit{Quantile Regression}} \\
QR stability (10 draws) & Most CV $< 8\%$ & Age CV = 78.6\% & Age unstable; all others sign-stable \\
Inter-quantile Wald & --- & 11/13 significant & Distribution varies for most attributes \\
\midrule
\multicolumn{4}{l}{\textit{Machine Learning}} \\
Random vs.\ geo holdout & $R^2 = 0.833$ & $R^2 = 0.425$ & 40.8 pp gap; spatial leakage \\
Remove geo features & Geo $R^2 = 0.425$ & Geo $R^2 = 0.519$ & $+$9.4 pp; geography hurts generalization \\
Add lat/lon & Random $R^2 = 0.870$ & Geo $R^2 = 0.464$ & Coordinates memorize location \\
Ablation: neighborhood & --- & $+$17.6 pp random & Largest single feature-group contribution \\
Ablation: full model & --- & $-$7.6 pp geo & Region dummies harm generalization \\
\midrule
\multicolumn{4}{l}{\textit{SHAP Stability}} \\
XGBoost vs.\ LightGBM & --- & $\rho = 0.990$ & Near-identical rankings \\
XGBoost vs.\ RF & --- & $\rho = 0.906$ & Consistent top-6 features \\
LightGBM vs.\ RF & --- & $\rho = 0.892$ & Consistent top-6 features \\
\bottomrule
\end{tabular}
\begin{tablenotes}
\footnotesize
\item \textit{Notes:} This table summarizes all robustness checks conducted in the paper. ``pp'' = percentage points. All findings refer to the same base sample of $N = 788{,}842$ unless noted. The geographic holdout $R^2$ of 0.425 for the full XGBoost model reflects the variant with region dummies included in the ablation design; the main specification yields 0.547 due to differences in training-set composition.
\end{tablenotes}
\end{threeparttable}
\end{table}
